\expandafter\ifx\csname CAPETINCLUDEMODE\endcsname\relax
\documentclass[11pt,a4paper]{article}
% preamble.tex — packages et commandes communes a toutes les fiches CAPET SII IE
\input{glyphtounicode}
\pdfgentounicode=1
\usepackage{cmap}
\usepackage[T1]{fontenc}
\usepackage[utf8]{inputenc}
\usepackage[french]{babel}
\usepackage[a4paper, top=1.6cm, bottom=1.6cm, left=1.4cm, right=1.4cm]{geometry}
\setlength{\headheight}{14.5pt}
\usepackage{amsmath, amssymb, amsfonts}
\usepackage{siunitx}
\sisetup{locale=FR, output-decimal-marker={,}, inter-unit-product=\ensuremath{{}\cdot{}}}
\usepackage{array, booktabs, multirow, tabularx}
\usepackage{xcolor}
\usepackage[most, breakable]{tcolorbox}
\usepackage{multicol}
\usepackage{enumitem}
\usepackage{mathtools}
\usepackage{lmodern}
\usepackage{microtype}
\usepackage{fancyhdr}
\usepackage{tikz}
\usetikzlibrary{positioning,arrows.meta}
\usepackage{lastpage}

% ── Couleurs par domaine ──────────────────────────────────────────────
\definecolor{colElec}{RGB}{26, 82, 148}       % Électrotechnique — bleu foncé
\definecolor{colElecLight}{RGB}{210, 225, 245} % fond bleu clair
\definecolor{colAuto}{RGB}{34, 120, 66}        % Automatique — vert
\definecolor{colMeca}{RGB}{185, 90, 20}        % Mécanique — orange
\definecolor{colTherm}{RGB}{160, 40, 40}       % Thermique — rouge
\definecolor{colQT}{RGB}{232, 244, 233}        % fond question type (vert très léger)
\definecolor{colJury}{RGB}{255, 243, 205}      % fond jury (jaune ambre)
\definecolor{colJuryFrame}{RGB}{200, 130, 0}   % bordure jury
\definecolor{colGrayLight}{RGB}{245, 245, 247} % fond gris neutre
\definecolor{colSeparator}{RGB}{180, 180, 180} % gris séparateur

% ── Style tcolorbox de section principale ────────────────────────────
\tcbset{
  ficheSection/.style={
    breakable,
    enhanced,
    colback=#1!7!white,
    colframe=#1,
    colbacktitle=#1,
    coltitle=white,
    fonttitle=\bfseries\normalsize,
    left=4pt, right=4pt, top=3pt, bottom=4pt,
    before skip=6pt, after skip=6pt,
  }
}

% ── Style encadré question type ──────────────────────────────────────
\tcbset{
  questiontype/.style={
    breakable,
    enhanced,
    colback=colQT,
    colframe=colElec!55!black,
    colbacktitle=colElec!55!black,
    coltitle=white,
    fonttitle=\bfseries\small,
    left=4pt, right=4pt, top=3pt, bottom=3pt,
    before skip=5pt, after skip=5pt,
  }
}

% ── Style encadré jury ───────────────────────────────────────────────
\tcbset{
  jury/.style={
    breakable,
    enhanced,
    colback=colJury,
    colframe=colJuryFrame,
    fonttitle=\bfseries\normalsize,
    left=4pt, right=4pt, top=3pt, bottom=3pt,
  }
}

% ── Mise en page ─────────────────────────────────────────────────────
\pagestyle{fancy}
\fancyhf{}
\renewcommand{\headrulewidth}{0.4pt}
\renewcommand{\footrulewidth}{0.3pt}

\setlist[itemize]{leftmargin=1.2em, itemsep=1pt, topsep=2pt, parsep=0pt}
\setlist[enumerate]{leftmargin=1.2em, itemsep=3pt, topsep=2pt, parsep=0pt}

\setlength{\parskip}{2pt}
\setlength{\columnsep}{10pt}

\def\CAPETBEGINFILE{\begin{document}}
\def\CAPETENDFILE{\end{document}}
\else
\def\CAPETBEGINFILE{}
\def\CAPETENDFILE{}
\fi

\fancyhead[L]{\textbf{\textcolor{colElec}{T23}} - Capteurs et Conditionnement du signal}
\fancyhead[R]{\small\textit{CAPET SII IE - Session 2026}}
\fancyfoot[C]{\small\thepage\,/\,\pageref{LastPage}}

\CAPETBEGINFILE

% -----------------------------------------------------------------------
%  BLOC EN-TeTE
% -----------------------------------------------------------------------

\begin{tcolorbox}[
  enhanced, breakable,
  colback=colElec,
  colframe=colElec,
  left=6pt, right=6pt, top=5pt, bottom=5pt,
  before skip=0pt, after skip=4pt
]
  {\color{white}\bfseries\large T23 \quad Capteurs et Conditionnement du signal}\\[2pt]
  {\color{colElec!30!white}\small
    \textbf{Domaine :} electronique - Instrumentation industrielle \quad
    \textbf{Presence :} \textbf{10/13} sujets (2013-2026) - theme \emph{incontournable}
  }\\[1pt]
  {\color{colElec!20!white}\footnotesize
    Sous-themes : PT100 - thermocouple - capteur inductif - jauge de contrainte - pont de Wheatstone - signal 4-20\,mA - montages AOP de conditionnement - mise - l'echelle CAN
  }
\end{tcolorbox}

% -----------------------------------------------------------------------
%  1. RAPPELS DE COURS
% -----------------------------------------------------------------------

\begin{tcolorbox}[
  ficheSection=colElec,
  title={1. Rappels de cours},
]

% -- 1.1 Taxonomie capteurs ---------------------------------------------
\textbf{1.1 Classification des capteurs}

\smallskip
\begin{center}
\renewcommand{\arraystretch}{1.25}
\begin{tabular}{lll}
\toprule
\textbf{Critere} & \textbf{Actif} & \textbf{Passif} \\
\midrule
Principe & Genere une FEM & Varie une impedance \\
Exemples & Thermocouple, piezoelectrique & PT100, jauge, potentiometre \\
Alimentation & Non (source autonome) & Oui (pont, courant) \\
\bottomrule
\end{tabular}
\end{center}

\medskip\hrule\vspace{4pt}

% -- 1.2 PT100 ---------------------------------------------------------
\textbf{1.2 Capteur PT100 (thermoresistance)}

\[
  R_T = R_0(1 + \alpha\,T)
  \qquad R_0 = 100\,\Omega \text{ - } 0\mathrm{C}
  \qquad \alpha = 3{,}85 \times 10^{-3}\,\mathrm{C}^{-1}
\]
\[
  \frac{\Delta R_T}{\Delta T} = R_0\,\alpha = 0{,}385\,\Omega/\mathrm{C}
\]
Capteur \textbf{passif} ; plage $-200^\circ$C - $+850^\circ$C. Necessite un pont d'alimentation.

\medskip\hrule\vspace{4pt}

% -- 1.3 Thermocouple --------------------------------------------------
\textbf{1.3 Thermocouple (effet Seebeck)}

Capteur \textbf{actif} : genere une FEM $e_{AB}$ de quelques mV entre deux metaux differents.
\begin{center}
\renewcommand{\arraystretch}{1.2}
\begin{tabular}{lll}
\toprule
\textbf{Type} & \textbf{Plage} & \textbf{Tolerance} \\
\midrule
J (Fe/Constantan) & $-210^\circ$C - $+1200^\circ$C & $\pm 2{,}5^\circ$C \\
K (Chromel/Alumel) & $-10^\circ$C - $+1100^\circ$C  & $\pm 2{,}5^\circ$C \\
\bottomrule
\end{tabular}
\end{center}
Necessite un transmetteur mV - 4-20\,mA ou un ampli d'instrumentation.

\medskip\hrule\vspace{4pt}

% -- 1.4 Pont de Wheatstone --------------------------------------------
\textbf{1.4 Pont de Wheatstone - jauge de contrainte}

\[
  V_{out} \approx \frac{V_{cc}}{4} \cdot \frac{\Delta R}{R_0} \quad (\Delta R \ll R_0)
\]
Avantage : compense les derives thermiques (4 jauges identiques). L'ampli differentiel amplifie $V_{out}$ jusqu'a la plage CAN.

\medskip\hrule\vspace{4pt}

% -- 1.5 Signal 4-20 mA ------------------------------------------------
\textbf{1.5 Signal normalise industriel 4-20\,mA}

\[
  x = x_{min} + \frac{I_s - 4}{16}\,(x_{max} - x_{min})
\]
\begin{itemize}
  \item $I_s = 4\,\mathrm{mA}$ - grandeur minimale (pas zero !)
  \item $I_s = 20\,\mathrm{mA}$ - grandeur maximale
  \item Immunite aux perturbations CEM sur grande longueur de ceble (2 fils)
\end{itemize}

\medskip\hrule\vspace{4pt}

% -- 1.6 Conditionnement AOP -------------------------------------------
\textbf{1.6 Montages AOP courants en conditionnement}

\begin{center}
\renewcommand{\arraystretch}{1.3}
\begin{tabular}{lll}
\toprule
\textbf{Montage} & \textbf{Gain} & \textbf{Usage} \\
\midrule
Inverseur & $-R_f/R_e$ & Amplification + inversion \\
Non-inverseur & $1 + R_f/R_e$ & Amplification sans inversion \\
Soustracteur (diff.) & $(R_2/R_1)(V^+ - V^-)$ & Pont de Wheatstone, PT100 \\
Suiveur & $1$ & Adaptation d'impedance \\
\bottomrule
\end{tabular}
\end{center}

\end{tcolorbox}

% -----------------------------------------------------------------------
%  2. FORMULAIRE SYNTHeTIQUE
% -----------------------------------------------------------------------

\begin{tcolorbox}[
  ficheSection=colElec,
  title={2. Formulaire synthetique},
]
\small
\begin{multicols}{2}
\renewcommand{\arraystretch}{1.45}
\begin{tabular}{@{}p{4.8cm}ll@{}}
\toprule
\textbf{Formule} & \textbf{Usage} & \textbf{U.} \\
\midrule
\multicolumn{3}{@{}l}{\footnotesize\textit{PT100}} \\
$R_T = R_0(1 + \alpha T)$ & Resistance & $\Omega$ \\
$\Delta R_T/\Delta T = R_0\alpha = 0{,}385$ & Sensibilite & $\Omega/C$ \\
\midrule
\multicolumn{3}{@{}l}{\footnotesize\textit{Pont Wheatstone}} \\
$V_{out} \approx (V_{cc}/4) \cdot (\Delta R/R_0)$ & Desequilibre & V \\
$V_{CAN} = K_a \cdot V_{out}$ & Apres ampli & V \\
\midrule
\multicolumn{3}{@{}l}{\footnotesize\textit{Signal 4-20\,mA}} \\
$x = x_{min} + (I_s-4)/16 \cdot (x_{max}-x_{min})$ & Mise - l'ech. & - \\
\bottomrule
\end{tabular}

\columnbreak

\renewcommand{\arraystretch}{1.45}
\begin{tabular}{@{}p{4.8cm}ll@{}}
\toprule
\textbf{Formule} & \textbf{Usage} & \textbf{U.} \\
\midrule
\multicolumn{3}{@{}l}{\footnotesize\textit{Conditionnement}} \\
$V_s = -(R_f/R_e) V_e$ & Inverseur & V \\
$V_s = (1+R_f/R_e) V_e$ & Non-inv. & V \\
$V_s = (R_2/R_1)(V^+-V^-)$ & Soustr. & V \\
\midrule
\multicolumn{3}{@{}l}{\footnotesize\textit{Capteur inductif}} \\
$f = N \cdot Z \cdot n / 60$ & Frequence signal & Hz \\
\midrule
\multicolumn{3}{@{}l}{\footnotesize\textit{Donnees PT100}} \\
$R_0 = 100\,\Omega$ ; $\alpha = 3{,}85\times10^{-3}\,C^{-1}$ & - & - \\
Plage : $-200^\circ$C - $+850^\circ$C & - & - \\
\bottomrule
\end{tabular}
\end{multicols}
\normalsize
\end{tcolorbox}

% -----------------------------------------------------------------------
%  3. QUESTIONS TYPES
% -----------------------------------------------------------------------

\begin{tcolorbox}[
  ficheSection=colElec,
  title={3. Questions types - extraites des sujets reels},
]

% -- QT1 -------------------------------------------------------------

% Mini-procedure de resolution type copie concours
\begin{tcolorbox}[questiontype,title={Procedure type copie concours (5 etapes)}]
\begin{enumerate}[leftmargin=1.5em]
  \item \textbf{Qualifier le modele} : hypotheses, domaine de validite, grandeurs d'entree/sortie.
  \item \textbf{Ecrire les lois} : equations de base avec symboles et unites SI.
  \item \textbf{Isoler l'inconnue} : expression litterale, puis application numerique tracee.
  \item \textbf{Controler l'ordre de grandeur} : coherence physique et unite finale.
  \item \textbf{Conclure concours} : phrase courte orientee decision technique.
\end{enumerate}
\end{tcolorbox}

% Visuel de methode : chemin de resolution CAPET
\begin{center}
\begin{tikzpicture}[font=\small, node distance=7mm and 6mm, >=Stealth]
  \node[draw, rounded corners, fill=colElec!8, inner sep=4pt] (n1) {Notions};
  \node[draw, rounded corners, fill=colElec!8, right=of n1, inner sep=4pt] (n2) {Methode};
  \node[draw, rounded corners, fill=colElec!8, right=of n2, inner sep=4pt] (n3) {Application annale};
  \node[draw, rounded corners, fill=colElec!8, right=of n3, inner sep=4pt] (n4) {Verification jury};
  \draw[->, thick, colElec] (n1) -- (n2);
  \draw[->, thick, colElec] (n2) -- (n3);
  \draw[->, thick, colElec] (n3) -- (n4);
\end{tikzpicture}
\end{center}
\begin{tcolorbox}[questiontype,
  title={QT1 - Identification du capteur et justification du conditionnement \hfill\normalfont\footnotesize E2 BAC3 2026 (Q1-Q2)}
]
\textbf{enonce.}
Capteur Jumo (reference 370) place dans une conduite. Identifier le type, le principe de mesure, sa nature active/passive. Justifier le besoin d'un conditionneur pour $T \in [-5\,;\,90^\circ\mathrm{C}]$, $V_{CAN} \in [0\,;\,3] \,\mathrm{V}$. Montrer que $c = 0{,}0316\,\mathrm{V/C}$.

\tcbline

\textbf{Corrige.}

Type : PT100 (thermoresistance) - transducteur resistif, capteur \textbf{passif}.

Plage capteur : variation de resistance correspondant - $95^\circ$C. Tension de sortie directe du capteur : trop faible et decalee - le conditionneur adapte la plage.

Coefficient : $c = \dfrac{V_{max} - V_{min}}{T_{max} - T_{min}} = \dfrac{3 - 0}{90 - (-5)} = \dfrac{3}{95} \approx \mathbf{0{,}0316\,\mathrm{V/C}}$
\end{tcolorbox}

\vspace{4pt}

% -- QT2 -------------------------------------------------------------

\begin{tcolorbox}[questiontype,
  title={QT2 - Pont de Wheatstone - expression de la tension de desequilibre \hfill\normalfont\footnotesize E2 Commune 2025}
]
\textbf{enonce.}
Un pont de Wheatstone alimente 4 jauges de contrainte ($R_0 = 120\,\Omega$, $V_{cc} = 1\,\mathrm{V}$). Une jauge subit $\Delta R = 0{,}0153\,\Omega$. (a) Expliquer l'interet du pont par rapport - une mesure directe. (b) Calculer $V_{out}$.

\tcbline

\textbf{Corrige.}

(a) Le pont compense les derives thermiques : si toutes les jauges varient de la meme faeon avec la temperature, la variation se cancelle. Seule la deformation mecanique appliquee - une jauge cree un desequilibre.

(b) Approximation petits desequilibres ($\Delta R \ll R_0$) :
\[
  V_{out} \approx \frac{V_{cc}}{4} \cdot \frac{\Delta R}{R_0} = \frac{1}{4} \cdot \frac{0{,}0153}{120} \approx \mathbf{31{,}9\,\mu\mathrm{V}}
\]
Gain ampli necessaire pour 5\,V en CAN : $K_a = 5 / (31{,}9 \times 10^{-6}) \approx 1{,}56 \times 10^5$.
\end{tcolorbox}

\vspace{4pt}

% -- QT3 -------------------------------------------------------------

\begin{tcolorbox}[questiontype,
  title={QT3 - Capteur inductif - calcul de frequence et choix \hfill\normalfont\footnotesize E1 2020}
]
\textbf{enonce.}
Capteur inductif detectant les dents d'un disque. Vitesse $n = 240\,\mathrm{tr/min}$, nombre de dents $Z = 24$. Calculer la frequence maximale du signal. Choisir le capteur adapte parmi la liste.

\tcbline

\textbf{Corrige.}

\[
  f = \frac{N \cdot Z \cdot n}{60} = \frac{1 \times 24 \times 240}{60} = \mathbf{96\,\mathrm{Hz}}
\]

Le capteur choisi doit avoir une frequence de commutation maximale $f_{max} \geq 96\,\mathrm{Hz}$. Justification chiffree : comparer $f_{max}$ du catalogue avec 96 Hz et selectionner le premier modele satisfaisant la condition.
\end{tcolorbox}

\end{tcolorbox}

% -----------------------------------------------------------------------
%  4. ATTENTES DU JURY
% -----------------------------------------------------------------------

\begin{tcolorbox}[
  jury,
  title={4. Attentes du jury - Erreurs recurrentes et points de vigilance},
]

\begin{enumerate}[label=\textbf{\arabic*.}, leftmargin=*, itemsep=5pt]

  \item \textbf{Identifier le capteur depuis la reference.} Lire le datasheet fourni ; ne pas deviner. La confusion PT100 (passif, resistance) / thermocouple (actif, FEM de l'ordre du mV) est frequente.

  \item \textbf{Justifier le choix chiffre.} Jury 2020 : "Le choix du capteur est rarement justifie." Toujours comparer avec les criteres numeriques (plage, tolerance, frequence).

  \item \textbf{Pont de Wheatstone.} Jury 2025 : moins de 20\,\% des candidats traitent correctement. Memoriser $V_{out} \approx (V_{cc}/4) \cdot (\Delta R/R_0)$ et verifier $\Delta R \ll R_0$.

  \item \textbf{Signal 4-20\,mA.} $4\,\mathrm{mA}$ correspond - la valeur minimale, pas - zero. La mise - l'echelle part de $I_s - 4$ et non de $I_s$.

  \item \textbf{Gain du conditionneur.} Calculer analytiquement : $G = (V_{max,CAN} - V_{min,CAN}) / (V_{max,capteur} - V_{min,capteur})$. Ne pas oublier l'offset si la plage capteur est decalee.

  \item \textbf{Relier capteur et CAN.} La tension de sortie du conditionneur doit \textbf{couvrir exactement} $[0\,;\,V_{ref}]$ - utiliser tout le codage du CAN pour maximiser la resolution effective.

\end{enumerate}

\medskip
\begin{tcolorbox}[
  colback=colJuryFrame!12!colJury,
  colframe=colJuryFrame,
  left=4pt, right=4pt, top=2pt, bottom=2pt,
  fontupper=\small\itshape
]
Rapport jury 2025 : "La partie capteurs a ete abordee par 85\,\% des candidats mais le pont de Wheatstone est correctement traite par moins de 20\,\%. Les candidats doivent calculer le gain de l'amplificateur differentiel en connaissant les limites des deux grandeurs."
\end{tcolorbox}

\end{tcolorbox}

% -----------------------------------------------------------------------
%  PIED - PReSENCE DANS LES SUJETS
% -----------------------------------------------------------------------

\begin{tcolorbox}[
  enhanced, breakable,
  colback=colGrayLight,
  colframe=colSeparator,
  left=4pt, right=4pt, top=3pt, bottom=3pt,
  before skip=4pt, after skip=0pt,
  fonttitle=\small\bfseries,
  title={Presence dans les sujets analyses (2013-2026)},
]
\begin{center}
\footnotesize\renewcommand{\arraystretch}{1.15}
\begin{tabular}{lccccccccccccc}
\toprule
\textbf{Annee} & 13 & 14 & 15 & 16 & 17 & 18 & 19 & 20 & 21 & 22 & 23 & 24 & 25 \\
\midrule
E1 (commune)   & - & $\checkmark$ & - & - & $\checkmark$ & - & - & $\checkmark$ & $\checkmark$ & $\checkmark$ & - & - & $\checkmark$ \\
E2 discipl. IE & - & - & - & - & - & - & $\checkmark$ & - & - & $\checkmark$ & $\checkmark$ & $\checkmark$ & $\checkmark$ \\
E2 appl. comm. & - & - & - & - & - & - & - & - & - & - & - & - & $\checkmark$ \\
\bottomrule
\end{tabular}
\end{center}
\smallskip
\footnotesize\textbf{Sous-themes frequents :}
PT100 ($R_T$, $\alpha$) - thermocouple (type J/K, FEM) - pont de Wheatstone - signal 4-20\,mA - montage AOP soustracteur - capteur inductif (frequence) - mise - l'echelle CAN

\medskip
\textbf{Ordres de grandeur (reperes rapides).}
\begin{itemize}[leftmargin=1.3em]
  \item Reseau BT FR : tension phase-neutre typique \SI{230}{\volt}, tension phase-phase typique \SI{400}{\volt}, frequence \SI{50}{\hertz}.
  \item Courants d'etude frequents en sujet : de quelques \si{\ampere} a quelques dizaines de \si{\ampere}.
  \item En conversion d'energie, un rendement attendu reste generalement entre 0{,}8 et 0{,}98 selon la technologie.
\end{itemize}

\medskip
\textbf{Synthese express (a refaire sans support).}
\begin{enumerate}[leftmargin=1.4em]
  \item Poser les hypotheses du modele et les conventions de signe.
  \item Ecrire 2 a 4 relations cles avec unites.
  \item Realiser un calcul numerique complet et verifier l'ordre de grandeur.
  \item Justifier le resultat par un argument technique de niveau CAPET.
\end{enumerate}

\end{tcolorbox}

\CAPETENDFILE









